\chapter{结果展示与示范分析}

\section{流程结果}

在模板工程层面，最重要的结果不是某一个统计量，而是结构能否稳定输出。当前项目已覆盖以下核心结果：

\begin{enumerate}
  \item 首页封面能够稳定输出分类号、UDC、密级、编号和双日期字段；
  \item 题名页能够集中展示博士后姓名、流动站、专业、单位和研究起止时间；
  \item 摘要、目录、缩略词表、正文、参考文献、附录和结尾材料均能进入统一编译链；
  \item 项目可同时支持完整仓库模式与单项目模式构建。
\end{enumerate}

\section{示范性结果描述}

若将本模板应用于真实课题，\figref{fig:postdoc-workflow} 所示的结构可以进一步展开为“数据来源说明、纳排标准、特征工程、模型解释和临床转化建议”等模块。当前公开示例仅保留最小但完整的写法骨架，以便后续使用者用自己的真实正文替换。

\section{模板验收重点}

对于博士后研究报告模板，人工验收时建议重点检查：

\begin{enumerate}
  \item 封面字段是否与学校或管理部门要求一致；
  \item 前置页顺序是否完整，页码是否连续；
  \item 正文标题层级、图表标题和参考文献是否统一；
  \item 个人简历、成果清单和永久通信地址等结尾部分是否遗漏。
\end{enumerate}
